\documentclass{article}
\usepackage[utf8]{inputenc}
\usepackage[T1]{fontenc}
\usepackage{geometry}
\geometry{
    a4paper,
    left=2.5cm,
    right=2.5cm,
    top=2.5cm,
    bottom=2.5cm
}
\usepackage{graphicx}
\usepackage{float}
\usepackage{hyperref}
\usepackage[english]{babel}
\usepackage{amsmath}
\usepackage{amssymb}
\usepackage{array}
\usepackage{booktabs}
\usepackage{caption}
\usepackage{subcaption}
\usepackage{enumitem}
\usepackage{fancyhdr}
\usepackage{titlesec}
\usepackage{url}
\usepackage{textcomp}
\usepackage{xcolor}
\usepackage{listings}
\usepackage{multirow}
\usepackage{longtable}
\usepackage{tabularx}
\usepackage{times}

% Set up page style
\pagestyle{fancy}
\fancyhf{}
\rfoot{\thepage}
\renewcommand{\headrulewidth}{0pt}

% Title formatting
\titleformat{\chapter}[display]{\normalfont\bfseries}{\chaptertitle\thechapter}{0pt}{\Huge}
\titlespacing*{\chapter}{0pt}{50pt}{40pt}

% Document info
\title{Template Document}
\author{Academic Support Agent}
\date{\today}

\begin{document}
\maketitle

\title{Modelo Relat rio CEFET}
	\CEFETopen
    \CEFETtitulo{Trabalho Legal:\\Esse trabalho ficou muito legal}


\section{Introdu o}
\lipsum[1-2]
\subsection{Coisas}
\lipsum[1-2]

%importando algoritmo de arquivo de codigo
\vspace{.3cm}
\lstinputlisting[language=C++]{code.cpp}

\subsection{Outras Coisas}
\lipsum[1-2]

%Escrevendo algoritmo
\begin{lstlisting}[language=Java, caption={},label={}]
package pessoa;

public class Pessoa {
    public String nome;
    public int idade;

    public static void main(String[] args) {
        Pessoa p1 = new Pessoa(); 
    
    }
}
\end{lstlisting}



\section{Coisas Dif ceis}
\lipsum[1]
\subsection{Coisas}

Isso   algo que \citet{Liu2007494} e \citet{law} disse em seu trabalho. A  gua   molhada \citep{1987element}. O sol   quente \citep{1987element,Liu2007494,site_sparselib}.

\lipsum[1-2]


%Tabela
\begin{center}
\vspace{.3cm}
\begin{tabular}{|l|l|l|}
\hline
\textbf{A}&\textbf{B}&\textbf{C}\\ \hline \hline
1&2&3\\
1&2&3\\
1&2&3\\ \hline
\end{tabular}
\vspace{.3cm}
\end{center}

\lipsum[1-2]

\subsection{Outras Coisas}
\lipsum[1-2]

De acordo com a Figura \ref{fig:figura2}.

\begin{figure}[h!]
\centering
\includegraphics[width=0.5\textwidth]{images/_logo}
\caption[]{Figura legal A}
\label{fig:figura1}
\end{figure}

\lipsum[1-2]

\begin{figure}[h]
  \centering
  \subfloat[]{\includegraphics[width=0.3\textwidth]{images/_logo}}\quad\quad
  \subfloat[]{\includegraphics[width=0.3\textwidth]{images/_logo}}
\caption[]{Figura legal B}
\label{fig:figura2}
\end{figure}

\lipsum[1-2]

\begin{eqnarray*}
(ano - nc) + (ano - ic) &\geq& 95\\
2\times ano - nc - ic  &\geq& 95\\
2\times ano &\geq& 95 + nc + ic\\
ano &\geq& (95 + nc + ic)/2
\end{eqnarray*}

\lipsum[1-2]
 
% Refer ncias-------------------

% alpha - Sorted alphabetically. Labels are formed from name of author and year of publication. 
% plain - Sorted alphabetically. Labels are numeric. 
% unsrt - Like plain, but entries are in order of citation. 
% abbrv - Like plain, but more compact labels. 

%
%

 %nome do arquivo bibtex

% -------------------------------
\end{document}